\documentclass[../main.tex]{subfiles}
\begin{document}

\section{Lý do chọn đề tài}
Trong thời đại chuyển đổi số, thị trường sách cũ tại Việt Nam vẫn chủ yếu diễn ra qua các hội nhóm mạng xã hội hoặc các cửa hàng truyền thống. Mô hình này gặp nhiều hạn chế:
\begin{itemize}
    \item \textbf{Rủi ro lừa đảo:} Người mua chuyển tiền nhưng không nhận được sách hoặc sách không đúng mô tả.
    \item \textbf{Khó khăn khi mua số lượng lớn:} Khó tìm kiếm và thanh toán cho nhiều loại sách từ các nguồn khác nhau.
    \item \textbf{Thiếu sự chuyên nghiệp:} Không có cơ chế đánh giá uy tín người bán và xử lý tranh chấp rõ ràng.
\end{itemize}
Vì vậy, việc xây dựng một nền tảng thương mại điện tử C2C dành riêng cho sách cũ với cơ chế thanh toán trung gian (Escrow) là vô cùng cần thiết.

\section{Mục tiêu của dự án}
Dự án tập trung vào việc phát triển nền tảng "Libris - C2C Bookstore" với các mục tiêu cụ thể:
\begin{itemize}
    \item Xây dựng hệ thống thanh toán trung gian an toàn, chỉ giải ngân cho người bán khi người mua đã hài lòng.
    \item Tối ưu hóa quy trình đặt hàng với tính năng tách đơn hàng tự động cho giỏ hàng đa người bán.
    \item Tích hợp cổng thanh toán điện tử VNPay để hiện đại hóa giao dịch.
    \item Xây dựng hệ thống ví điện tử nội bộ để quản lý dòng tiền và đối soát hoa hồng.
\end{itemize}

\section{Đối tượng và phạm vi nghiên cứu}
\begin{itemize}
    \item \textbf{Đối tượng:} Các cá nhân có nhu cầu thanh lý sách cũ hoặc tìm mua sách với chi phí tiết kiệm.
    \item \textbf{Phạm vi:} Tập trung vào các giao dịch sách tại thị trường Việt Nam, sử dụng các công nghệ Web hiện đại nhất hiện nay.
\end{itemize}

\end{document}
